\subsection{Unpacking the identity}
\label{sec:identity}

The reason this leads directly to \rs{} is an accounting identity, and it is worth writing
out in full because the interesting term is the one usually waved away.

Start from the definition of a long yield as an average of expected short rates plus a
residual:
\begin{equation}
y^{(n)}_t \;=\; \underbrace{\frac{1}{n}\sum_{j=0}^{n-1} E_t\!\left[i_{t+j}\right]}_{\text{expectations component } rn^{(n)}_t} \;+\; TP^{(n)}_t .
\label{eq:eh}
\end{equation}
This is a definition, not an assumption: $TP^{(n)}_t$ is whatever is left over.

Now split the nominal policy rate into a neutral level and a deviation from it. Define the
\emph{neutral nominal rate} and the \emph{policy gap}
\begin{equation}
i^{*}_t \;\equiv\; r^{*}_t + \pi^{*}_t ,
\qquad
\text{gap}_t \;\equiv\; i_t - i^{*}_t \;=\; \underbrace{\left(i_t - E_t\pi_{t+1} - r^{*}_t\right)}_{\text{real policy stance}} \;+\; \underbrace{\left(E_t\pi_{t+1} - \pi^{*}_t\right)}_{\text{inflation deviation}} ,
\label{eq:gapdef}
\end{equation}
where $r^{*}_t$ is the equilibrium real rate and $\pi^{*}_t$ the long-run inflation
anchor. The gap is therefore \textbf{the stance of policy}: how far the Fed is holding the
short rate above or below the level that would be neutral, plus how far inflation is
expected to run from its anchor. Under a Taylor-type rule
$i_t = i^{*}_t + \phi_\pi(\pi_t-\pi^{*}_t) + \phi_x x_t + v_t$, the gap is exactly the
systematic and discretionary response to the cycle,
\begin{equation}
\text{gap}_t \;=\; \phi_\pi\!\left(\pi_t-\pi^{*}_t\right) + \phi_x x_t + v_t .
\label{eq:taylor}
\end{equation}
It is a cyclical, mean-reverting object: it is large in 1994, in 2008 and in 2023, and zero
on average.

Substituting \eqref{eq:gapdef} into \eqref{eq:eh} and adding the assumption that the
endpoints are (locally) martingales, $E_t[r^{*}_{t+j}]=r^{*}_t$ and
$E_t[\pi^{*}_{t+j}]=\pi^{*}_t$, gives the identity used throughout:
\begin{equation}
y^{(n)}_t \;=\; r^{*}_t \;+\; \pi^{*}_t \;+\; \frac{1}{n}\sum_{j=0}^{n-1} E_t\!\left[\text{gap}_{t+j}\right] \;+\; TP^{(n)}_t .
\label{eq:identity}
\end{equation}

\subsubsection*{How much does the gap term actually contribute?}

The usual move at this point is to say the gap ``averages out'' at a ten-year horizon. That
is too quick, and it is the step I initially got wrong. If the gap mean-reverts as an AR(1)
with monthly persistence $\rho$, the expected-gap term collapses to a single scalar times
today's stance,
\begin{equation}
\frac{1}{n}\sum_{j=0}^{n-1} E_t\!\left[\text{gap}_{t+j}\right]
\;=\; \omega(\rho,n)\,\text{gap}_t ,
\qquad
\omega(\rho,n) \;=\; \frac{1-\rho^{\,n}}{n\,(1-\rho)} ,
\label{eq:omega}
\end{equation}
and for a forward rate covering periods $n_1$ through $n_1+n_2-1$ the analogous weight is
$\omega^{f} = \rho^{n_1}\left(1-\rho^{n_2}\right)\big/\left[n_2(1-\rho)\right]$.

Two objects enter $\omega$: the maturity $n$, which is known, and the persistence $\rho$,
which is a \emph{belief} and is not directly observed. So the relevant exercise is a
sensitivity analysis over $\rho$, expressed as the half-life of a stance deviation --- the
time markets expect a 100 bp deviation from neutral to decay to 50 bp:

\begin{center}\small
\begin{tabular}{lcccc}
\toprule
Half-life of the gap & $\omega$, 2-year & $\omega$, 5-year & $\omega$, 10-year & $\omega^{f}$, 5y5y forward\\ \midrule
1 year  & 0.56 & 0.29 & 0.15 & 0.01\\
2 years & 0.73 & 0.48 & 0.28 & 0.09\\
3 years & 0.81 & 0.60 & 0.39 & 0.19\\
4 years & 0.85 & 0.67 & 0.48 & 0.28\\
5 years & 0.88 & 0.73 & 0.54 & 0.36\\
\bottomrule\end{tabular}
\end{center}

\noindent The upper rows are not the empirically relevant ones. Estimating an AR(1) on a
crude gap proxy --- the one-year yield less (HLW $r^{*}$ plus 2) --- at monthly frequency gives
a half-life of 4.4 years on the full sample, 5.0 years post-1985 and 4.1 years post-1995. The
bottom rows of the table are where the data sits, so $\omega(10\text{y}) \approx 0.5$: a 100 bp
revision to the expected stance --- with no change whatsoever in \rs{}, $\pi^{*}$ or the term
premium --- moves the ten-year yield by roughly half as much.

That estimate is an upper bound and should be treated as one. Any slow drift in the true
neutral rate that the proxy misses --- the disinflation, movements in $r^{*}$ that HLW does not
capture --- is dumped into the residual and reads as gap persistence. A better $\pi^{*}$ series
would pull the half-life down. But it would have to fall to roughly one year before
$\omega(10\text{y})$ reached 0.15.

\subsubsection*{Why this bears on the far-forward evidence}

The last column is the one that matters for Hillenbrand's second piece of evidence
(\S1.2). Testing the pattern on the 5-year, 5-year-forward rate is exactly the right
instinct: the first five years of the expected path drop out, so the stance term is
mechanically damped. But the test is only \emph{decisive} if the gap is fast. At a one-year
half-life the 5y5y weight is 0.01 and the argument is airtight; at four years it is 0.28,
and a 100 bp stance revision still moves the 5y5y forward by nearly 30 bp --- comfortably
enough to generate the observed window moves without any \rs{} news at all.

At the persistence actually estimated above, the far-forward result does not separate
endpoint news from stance news. Whether it does is an empirical question about $\rho$ that
this literature has not addressed, and it is cheap to address: the half-life of the policy
gap is estimable, and the sensitivity of the conclusion to it is one line of arithmetic.

\subsubsection*{The identification problem, stated}

Differencing \eqref{eq:identity} over an announcement window,
\begin{equation}
\Delta y^{(n)}_t \;=\; \underbrace{\Delta r^{*}_t + \Delta \pi^{*}_t}_{\text{endpoint news}}
\;+\; \underbrace{\omega\,\Delta\text{gap}_t + \text{gap}_t\,\Delta\omega}_{\text{stance and persistence news}}
\;+\; \Delta TP^{(n)}_t .
\label{eq:windowdiff}
\end{equation}
Four terms, one observable. Note that the second stance term, $\text{gap}_t\,\Delta\omega$, is a
revision to \emph{how long the current stance will last} --- which is what forward guidance is,
and which is largest precisely when the stance is large.

\subsubsection*{Three consequences for the project}

\textbf{1. ``The Fed moves long rates'' does not imply ``the Fed moves \rs{} beliefs.''} The
stance terms are a live alternative at the persistence the data suggest, and the
$\text{gap}_t\,\Delta\omega$ piece is largest exactly when the stance is large --- 2008--09 and
2022--24. It is also plausibly what Bauer, Pflueger \& Sunderam's perceived-reaction-function
result is picking up.

\textbf{2. $\pi^{*}$ is not fixed over most of the sample.} The claim that the inflation
anchor is constant is defensible after roughly 2000 and indefensible before it: the 1989--1999
window declines coincide with the disinflation and with the Fed acquiring credibility, so
$\Delta\pi^{*}$ plausibly carries a large share of them. Since 81\% of the in-window decline
predates 2012 (\S5.3), this is not a small caveat --- it means the pre-2012 evidence cannot
be read as \rs{} learning either. TIPS breakevens begin in 1997 (liquid from about 2003), so
the decomposition is only testable on the back half.

\textbf{3. The identity is what makes the term-premium question sharp.} If $\Delta\pi^{*}$ and
the stance terms could be shut down, then a permanent window move in the ten-year would have
to be $\Delta r^{*}$ or $\Delta TP$, and the whole question reduces to splitting those two.
That is the design in \S5 --- and it is why the failure of the market-side endpoint measure
is fatal rather than inconvenient.
